% LaTeX Template für Abgaben an der Universität Stuttgart
% Autor: Sandro Speth
% Bei Fragen: Sandro.Speth@iste.uni-stuttgart.de
%-----------------------------------------------------------
% Hauptmodul des Templates: Hier können andere Dateien eingebunden werden
% oder Inhalte direkt rein geschrieben werden.
% Kompiliere dieses Modul um eine PDF zu erzeugen.

% Dokumentenart. Ersetze 12pt, falls die Schriftgröße anzupassen ist.
\documentclass[12pt]{scrartcl}
% Einbinden der Pakete, des Headers und der Formatierung.
% Mit den \include und \input Befehlen können Dateien eingebunden werden:
% \include: Fügt einen Seitenumbruch nach dem Text ein
% \input: Fügt KEINEN Seitenumbruch nach dem Text ein
% LaTeX Template für Abgaben an der Universität Stuttgart
% Autor: Sandro Speth
% Forked Version: JSendra
%-----------------------------------------------------------
% Modul fuer verwendete Pakete.
% Neue Pakete einfach einfuegen mit dem \usepackage Befehl:
% \usepackage[options]{packagename}

% Encoding and Fonts
\usepackage[utf8]{inputenc}
\usepackage[T1]{fontenc}
\usepackage[ngerman]{babel}
\usepackage{lmodern}

% Graphics and Colors
\usepackage{graphicx}
\usepackage[pdftex,hyperref,dvipsnames,table]{xcolor}
\usepackage{tikz}
\usetikzlibrary{automata,positioning,arrows.meta}

% Math packages
\usepackage{amsmath,amssymb,amstext,amsthm}

% Lists, URLs, etc.
\usepackage[inline]{enumitem} % Ermöglicht ändern der enum Item Zahlen
\usepackage{hyperref}
\usepackage{url}
\usepackage{booktabs}
\usepackage{float}
\usepackage{pifont}
\usepackage{tcolorbox}

% Plotting
\usepackage{pgfplots}

% Code Listings
\usepackage{listings}
\usepackage{minted}

% UML Class Diagrams
% One Day I may fix the UML Diagrams
%\usepackage{styles/tikz-uml}
%\usetikzlibrary{positioning}

% Algorithms
\usepackage[lined,algonl,boxed]{algorithm2e}
% alternative zu algorithm2e:
%\usepackage[]{algorithm} %counter mit chapter
%\usepackage{algpseudocode}

% Geometry
\usepackage[a4paper,lmargin={2cm},rmargin={2cm},tmargin={3.5cm},bmargin = {2.5cm},headheight = {4cm}]{geometry}

% Headers
\usepackage[headsepline]{scrlayer-scrpage} 
\pagestyle{scrheadings} 

\pgfplotsset{compat=1.18}


% LaTeX Template für Abgaben an der Universität Stuttgart
% Autor: Sandro Speth
% Bei Fragen: Sandro.Speth@iste.uni-stuttgart.de
% Contributor: JSendra - extended and modified for current Abgaben schema
%-----------------------------------------------------------
% Modul beinhaltet Befehl fuer Aufgabennummerierung,
% sowie die Header Informationen.

% Überschreibt enumerate Befehl, sodass 1. Ebene Items mit
\renewcommand{\theenumi}{(\alph{enumi})}
% (a), (b), etc. nummeriert werden.
\renewcommand{\labelenumi}{\text{\theenumi}}

% Counter für das Blatt und die Aufgabennummer.
% Ersetze die Nummer des Übungsblattes und die Nummer der Aufgabe
% den Anforderungen entsprechend.
% Gesetz werden die counter in der hauptdatei, damit siese hier nicht jedes mal verändert werden muss
% Beachte:
% \setcounter{countername}{number}: Legt den Wert des Counters fest
% \stepcounter{countername}: Erhöht den Wert des Counters um 1.
\newcounter{sheetnr}
\newcounter{exnum}

\newcounter{teilaufgabe}[exnum]


% Befehl für die Aufgabentitel
\newcommand{\exercise}[1]{
    \section*{
        Aufgabe \theexnum\stepcounter{exnum}: #1
        }
} % Befehl für Aufgabentitel

\newcommand{\teil}[1]{
    \section*{
        Aufgabe \theexnum.\theteilaufgabe\stepcounter{teilaufgabe}: #1
    }
}


% Formatierung der Kopfzeile
% \ohead: Setzt rechten Teil der Kopfzeile mit
% Namen und Matrikelnummern aller Bearbeiter
\ohead{Josse Sendra (3542251)\\
Dennis Hehn (3593682)\\
Matthias Wagner (3542811)}
% \chead{} kann mittleren Kopfzeilen Teil sezten
% \ihead: Setzt linken Teil der Kopfzeile mit
% Modulnamen, Semester und Übungsblattnummer
\ihead{NumStatStoch\\
Sommersemester 26\\
Blatt \thesheetnr}




% Counter für das Blatt und die Aufgabennummer.
% Ersetze die Nummer des Übungsblattes und die Nummer der Aufgabe
% den Anforderungen entsprechend.
% Definiert werden die Counter in FormatAndHeader.tex
% Beachte:
% \setcounter{countername}{number}: Legt den Wert des Counters fest
% \stepcounter{countername}: Erhöht den Wert des Counters um 1.
\setcounter{sheetnr}{2} % Nummer des Übungsblattes
\setcounter{exnum}{1} % Nummer der Aufgabe

% Beginn des eigentlichen Dokuments
\begin{document}
% Nutze den \exercise{Aufgabenname} Befehl, um eine neue Aufgabe zu beginnen.
% Möchtest du eine Aufgabe in der Nummerierung überspringen, schreibe vor der Aufgabe: \stepcounter{exnum}
% Möchtest du die Nummer einer Aufgabe auf eine beliebige Zahl x setzen, schreibe vor der Aufgabe: \setcounter{exnum}{x}

% Aufgabe 1
\exercise{Programmiersprachen}

\begin{enumerate}
    \item Diskutieren Sie den Unterschied zwischen kompilierten und interpretierten Programmiersprachen.

        \textbf{Erläuterung:}

        \begin{itemize}
            \item \textbf{Kompiliert:} Der Quellcode wird vor der Ausführung vollständig in Maschinensprache übersetzt (z.\,B. \texttt{C}, \texttt{C++}). Dadurch läuft das Programm schnell, ist aber oft an ein bestimmtes Betriebssystem oder eine Architektur gebunden.
    
            \item \textbf{Interpretiert:} Der Quellcode wird Zeile für Zeile oder Anweisung für Anweisung zur Laufzeit vom Interpreter ausgeführt (z.\,B. \texttt{Python}, \texttt{JavaScript}). Das erleichtert Testen und Entwicklung, aber die Ausführung ist langsamer, da die Übersetzung ständig stattfindet.
        \end{itemize}

   

    \item Nennen Sie je einen Vorteil und einen Nachteil einer kompilierten bzw.
einer interpretierten Sprache.

    \begin{table}[H]
    \centering
    \begin{tabular}{@{}l|l|l@{}}
    \toprule
        Typ & Vorteil & Nachteil \\ \midrule
        Kompiliert & Sehr effizient und schnell & Plattformabhängig, längere Übersetzungszeit \\
        \\
        Interpretiert & Einfach zu testen und zu debuggen, & Langsamere Ausführung,\\
                      & plattformunabhängig & benötigt Interpreter zur Laufzeit \\
    \bottomrule
    \end{tabular}
    \end{table}
    
\end{enumerate}

% Aufgabe 2
\exercise{Datentypen}
% Teilaufgaben
\begin{enumerate}
	\item Nennen Sie den Datentyp jeder Variable im folgenden Code-Ausschnitt.
        
        \lstdefinestyle{javaStyle}{
            language=Java,
            backgroundcolor=\color{gray!5},
            basicstyle=\ttfamily\small,
            keywordstyle=\color{blue!80!black}\bfseries,
            stringstyle=\color{teal!70!black},
            commentstyle=\color{gray!80}\itshape,
            numbers=left,
            numberstyle=\tiny\color{gray},
            stepnumber=1,
            frame=single,
            framerule=0.5pt,
            breaklines=true,
            showstringspaces=false,
            tabsize=2,
            morekeywords={var},
        }

    \
    \begin{flushleft}
       \begin{lstlisting}[style=javaStyle, caption={Code-Ausschnitt zur Bestimmung der Datentypen}, xleftmargin=0pt, framexleftmargin=0pt]
var variable1 = 10 / 3;
var variable2 = "10";
var variable3 = 15f;
var variable4 = "15" + 10;
var variable5 = 15.0;
var variable6 = 15.0d;
var variable7 = variable1 + variable4;
var variable8 = variable5 + variable6;
var variable9 = (int)(double) variable5;
var variable10 = variable3 <= variable5;
        \end{lstlisting}
        \end{flushleft}        
        \begin{table}[H]
        \centering
        \begin{tabular}{@{}l|l|p{8cm}@{}}
        \toprule
            \textbf{Variable} & \textbf{Datentyp} & \textbf{Erklärung} \\ \midrule
            variable1 & \texttt{int} & \texttt{10 / 3} $\rightarrow$ Ganzzahlige Division $\rightarrow$ Ergebnis \texttt{3} \\
            variable2 & \texttt{String} & In Anführungszeichen $\rightarrow$ Text \\
            variable3 & \texttt{float} & \texttt{15f} $\rightarrow$ \texttt{f} kennzeichnet Literal als \texttt{float} \\
            variable4 & \texttt{String} & \texttt{"15" + 10} $\rightarrow$ String-Verkettung $\rightarrow$ \texttt{"1510"} \\
            variable5 & \texttt{double} & \texttt{15.0} $\rightarrow$ Fließkommazahl, standardmäßig \texttt{double} \\
            variable6 & \texttt{double} & \texttt{15.0d} $\rightarrow$ explizit als \texttt{double} markiert \\
            variable7 & \texttt{String} & \texttt{variable1 + variable4} $\rightarrow$ \texttt{int + String} $\rightarrow$ automatische Umwandlung zu \texttt{String} \\
            variable8 & \texttt{double} & \texttt{variable5 + variable6} $\rightarrow$ \texttt{double + double} $\rightarrow$ Ergebnis \texttt{double} \\
            variable9 & \texttt{int} & \texttt{(int)(double) variable5} $\rightarrow$ zuerst \texttt{double}-Cast, dann \texttt{int}-Cast $\rightarrow$ \texttt{int} \\
            variable10 & \texttt{boolean} & \texttt{variable3 <= variable5} $\rightarrow$ Vergleich $\rightarrow$ \texttt{true} oder \texttt{false} \\ \bottomrule
        \end{tabular}
        \end{table}
    
	\item Widening Conversions
    
    %\begin{figure}[!h]
    %\centering
    %\begin{tikzpicture}[roundnode/.style={circle, draw=black!60,very thick, minimum size=7mm}]

    %Nodes
    %\node[roundnode] (first) {A};
    %\node[roundnode] (second) [right=of first]{B};
    %\node[roundnode] (third) [right=of second]{C};

    %Lines
   % \draw[->] (first.east) -- (second.west);
   % \draw[->] (second.east) -- (third.west);

    \textbf{Widening Conversions (Java Primitive Types)}

\textbf{Definition:} Eine Widening Conversion ist eine automatische Typkonvertierung in Java, bei der ein kleinerer primitiver Datentyp in einen größeren Typ konvertiert wird, ohne Datenverlust.

\textbf{Graphische Darstellung:}


\begin{tikzpicture}[node distance=1cm, >=Stealth, every node/.style={draw, rounded corners, minimum width=1.5cm, align=center}]
    \node (byte) {byte};
    \node (short) [right=of byte] {short};
    \node (int) [right=of short] {int};
    \node (long) [right=of int] {long};
    \node (float) [right=of long] {float};
    \node (double) [right=of float] {double};
    \node (char) [right=of double] {char};

    % Arrows
    \draw[->] (byte) -- (short);
    \draw[->] (short) -- (int);
    \draw[->] (int) -- (long);
    \draw[->] (long) -- (float);
    \draw[->] (float) -- (double);

    \draw[->] (double) -- (char);
   % \draw[->] (int) -- (long); % optional, already in transitives
\end{tikzpicture}


\textbf{Erläuterung:}  
- Kleine Datentypen in Java werden automatisch in größere Typen konvertiert. \\ 
- Beispiel: \texttt{int -> double} ist gültig; \texttt{double -> int} ist nicht automatisch möglich.  \\
- Transitiv abgeleitete Konvertierungen (z.\,B. \texttt{byte -> int}) müssen nicht extra gezeichnet werden.
    
    
\end{enumerate}
\newpage
%Aufgabe 3
\exercise{Boolesche Logik}

\begin{enumerate}

    %a)
    \item Was ist die Ausgabe der Methode example1?

        \begin{lstlisting}[style=javaStyle]
public static void example1() {
    boolean a = true;
    boolean b = false;
    int x = 10;
    int y = 20;
    int z = 30;
        
    if (a || b) {
        System.out.println("C1");
        
            if (x < y && x < z) {
                System.out.println("C2");
            } else if (x < y || x < z) {
                System.out.println("C3");
            } else {
                System.out.println("C4");
            }
        }
}
\end{lstlisting}

    \textbf{Schritt-für-Schritt Analyse:}
\begin{itemize}
    \item Bedingung \texttt{a || b} $\Rightarrow$ \texttt{true || false} $\Rightarrow$ \texttt{true} $\Rightarrow$ Ausgabe: \texttt{C1}
    \item Bedingung \texttt{x < y && x < z} $\Rightarrow$ \texttt{10 < 20 && 10 < 30} $\Rightarrow$ \texttt{true} $\Rightarrow$ Ausgabe: \texttt{C2}
\end{itemize}

\textbf{Ausgabe:}
\begin{verbatim}
C1
C2
\end{verbatim}

    %b)
    \newpage
    \item Was ist die Ausgabe der Methode example2?

    \begin{lstlisting}[style=javaStyle]
public static void example2() {
    int x = 10;
    int y = 20;
    int z = 30;
    int i = 0;

    while (i < 10) {
        if (x / 10 > i || i * x > z) {
            System.out.println("L1");
            i++;
        } else if (i > y / 3) {
            System.out.println("L2");
            i += 2;
        } else {
            System.out.println("L3");
            i += 3;
        }

        if (i % 2 == 0) {
            System.out.println("Even: " + i);
            i++;
        }
    }
}
\end{lstlisting}

\textbf{Schritt-für-Schritt Analyse:}

\begin{tabular}{@{}c|c|c|c@{}}
\toprule
i & Bedingung & Ausgabe & Neuer i \\ \midrule
0 & 10/10>0 $\Rightarrow$ 1>0 true & L1 & 1 \\
1 & 1>1 false, 1*10>30 false $\Rightarrow$ else & L3, Even:4 & 5 \\
5 & 1>5 false, 5*10>30 true & L1, Even:6 & 7 \\
7 & 1>7 false, 7*10>30 true & L1, Even:8 & 9 \\
9 & 1>9 false, 9*10>30 true & L1 & 10 (Ende) \\ \bottomrule
\end{tabular}

\textbf{Ausgabe:}
\begin{verbatim}
L1
L3
Even: 4
L1
Even: 6
L1
Even: 8
L1
\end{verbatim}

    %c)
    \newpage
    \item Geben Sie die Warheitstabellen für die folgenden booleschen Ausdrücke.

    1. \texttt{(!A && B) || (C && !D) || (C && B)}  

\begin{tabular}{@{}c|c|c|c|c@{}}
\toprule
A & B & C & D & Ergebnis \\ \midrule
F & F & F & F & F \\
F & F & F & T & F \\
F & F & T & F & T \\
F & F & T & T & T \\
F & T & F & F & T \\
F & T & F & T & T \\
F & T & T & F & T \\
F & T & T & T & T \\
T & F & F & F & F \\
... & ... & ... & ... & ... \\ \bottomrule
\end{tabular}\\

2. \texttt{!(A || B || C || D)}\\
\begin{tabular}{@{}c|c|c|c|c@{}}
\toprule
A & B & C & D & Ergebnis \\ \midrule
F & F & F & F & T \\
\text{else combinations} & & & & F \\ \bottomrule
\end{tabular}\\

3. \texttt{!A \&\& !B \&\& !C \&\& !D}\\
\begin{tabular}{@{}c|c|c|c|c@{}}
\toprule
A & B & C & D & Ergebnis \\ \midrule
F & F & F & F & T \\
\text{else combinations} & & & & F \\ \bottomrule
\end{tabular}\\

\end{enumerate}

% Ende des Dokuments
\end{document}
